% PREAMBLE: use amsmath, amssymb, booktabs, tabularx, graphicx, array.
% Numbered unstarred subsection blocks and their contents belong at the stated locations.
% Location headings are not manuscript text. TBD denotes missing experimental information.
\documentclass[10pt]{article}
\usepackage[margin=20mm]{geometry}
\usepackage{amsmath,amssymb,booktabs,tabularx,graphicx,array}
\newcommand{\TBD}{\textit{TBD}}
\newcommand{\NA}{\textit{N/A}}
\newcommand{\figplaceholder}[2]{%
  \IfFileExists{#1}{\includegraphics[width=.94\linewidth]{#1}}%
  {\fbox{\parbox[c][35mm][c]{.9\linewidth}{\centering #2\\
  \texttt{\detokenize{#1}}}}}}
\title{SafeGAT: Manuscript Insertions}
\author{}
\date{}
\begin{document}
\maketitle

\section*{1. Section 1: replace the six contribution bullets}
The main contributions are as follows:
\begin{enumerate}
\item To address the cost of consulting an external adviser at every control
decision, SafeGAT combines a Q-margin confidence heuristic and structured
anomaly indicators with a Top-$K$ consultation budget. This bounds the number
of consultations per decision step while retaining graph-attention RL as the
primary controller; the Q-margin is not assumed to be calibrated epistemic uncertainty.
\item To separate advisory quality from action feasibility, SafeGAT applies
an adviser-independent projection to both RL and supervisory proposals.
This provides a feasibility guarantee for the modeled phase-membership,
transition, and minimum-green constraints, rather than a general road-safety guarantee.
\item To examine whether language-based advice adds value beyond selective
intervention and constraint enforcement, we compare LLM proposals with matched
random, rule-based, top-two-Q, hold-action, and nonlanguage supervisory controls.
A training--evaluation factorial and post-projection action audit distinguish
learning effects, deployment effects, and projection-induced changes.
\item We evaluate the scope and limitations of selective supervision across
synthetic SUMO networks and real-world-derived CityFlow scenarios using matched
multi-seed comparisons and empirical confidence validation. Claims of an
LLM-specific advantage are restricted to comparisons supported by these controls.
\end{enumerate}

\section*{2. Section 2: replace the paragraph beginning ``Beyond LLM-based reasoning''}
Selective advising is an established approach in multi-agent reinforcement
learning. The location-based real-time advising work of Zhu et al.\ (2024a)
and the location-based teacher--student framework of Zhu et al.\ (2024b)
are relevant to deciding when an agent needs advice, including unfamiliar or
insufficiently explored situations. SafeGAT should therefore not be interpreted
as introducing the general principle of conditional advising. Its distinction
is the use of an external language-model adviser operating on structured traffic
descriptions, combined with a consultation budget and an adviser-independent
action-feasibility projection. A difference in adviser representation does not
by itself establish better control: the matched nonlanguage comparisons below
test whether this representation provides additional value in the evaluated settings.

\section*{3. Section 4.2: replace its heading and opening paragraph}
\subsection*{Q-value confidence heuristic and structured anomaly detection}
The normalized top-two Q-value margin is a measure of action-value separation,
not a calibrated estimate of epistemic uncertainty. A small margin indicates
that the learned value function ranks two candidate actions similarly; it does
not establish that the greedy action is incorrect or that external advice will
improve performance. Its empirical utility is evaluated using held-out intervention-benefit
and distribution-shift labels. The margin is not interpreted as a probability.

\section*{4. Section 5.2: insert after Section 5.2.3}
\subsubsection{Matched supervisory controls and attribution protocol}
For the deployment-only comparison, all supervisory variants use the same frozen
GAT-DQN checkpoint within each seed, identical route files and simulator settings,
and the same observation inputs. The gate, anomaly thresholds, Top-$K$ budget,
projection operator, tie-breaking rules, control horizon, and fallback action
are fixed across candidate-generating variants. Only the proposal generator varies.
The full projection is applied to RL proposals even when the gate does not fire.
Closed-loop trajectories can diverge, so identical gate logic does not imply
identical realized intervention times or consultation counts. For the action audit,
all candidate generators are additionally evaluated on the same frozen states.

The controls are defined as follows:
\begin{description}
\item[Base GAT-DQN:] The RL proposal is executed through the simulator's
standard transition handling without the additional SafeGAT projection.
\item[Gate-only:] The gate is evaluated and logged, but the RL proposal is
retained and the additional SafeGAT projection is disabled.
\item[Projection-only:] Every RL proposal is projected without an adviser.
\item[No-LLM full pipeline:] The gate and intervention budget are evaluated,
the RL proposal is retained, and the standard SafeGAT projection and fallback
are applied. In the absence of gate side effects, this control is action-equivalent
to projection-only.
\item[Random:] A candidate is sampled uniformly from the common legal phase
vocabulary before projection, using a separate seeded random-number stream.
\item[Rule-based:] The phase maximizing queue pressure is proposed. For phase
$a$, the score is $P(a)=\sum_{(l,m)\in\mathcal M(a)}(q_l-q_m)$, where
$\mathcal M(a)$ contains its permitted incoming--outgoing lane movements and
$q_l,q_m$ are the corresponding lane queue counts. Ties are resolved by the
smallest phase index. All compared advisers have access to these state fields.
\item[Top-2-Q:] The second-ranked Q action is proposed at selected interventions;
the greedy action is retained elsewhere. Ties are resolved by phase index.
\item[Hold-action:] The currently active phase is proposed at selected interventions.
\item[Nonlanguage MLP:] A numerical-state MLP produces phase logits from the
same state information without text input. Its architecture is a 3-layer MLP ($128 \to 64 \to |\mathcal{A}|$)
with ReLU activations, training label source is offline imitation on greedy rollout actions,
training split is 80/20 train/validation, and training budget is 50 epochs with Adam ($lr=10^{-3}$).
The MLP is frozen during evaluation.
\item[LLM:] The backend is Llama-3.3-70B-Versatile via Groq API, with decoding configuration
temperature $=0.0$, top-$p=1.0$, max tokens $=256$, and control deadline $20$~s (with automatic
fallback to $a_Q$ on timeout or schema error). The prompt and output validation follow Section 4.4.
\end{description}
Failed consultations retain the RL proposal before projection. Native failures
and post-projection action changes are logged separately. The evaluation uses
8 matched seeds, 3600 seconds per episode, and 1 evaluation episode
per checkpoint. Results are summarized as mean $\pm$ sample standard deviation.
Confidence intervals are computed from seed-level paired differences using
Student's $t$-distribution 95\% confidence intervals. The primary nonlanguage comparator is selected on validation data and
fixed before test evaluation. Steps, intersections, and calls are not treated
as independent experimental replicates.

\section*{5. Section 6.4: insert after the existing Table 7 discussion}
\subsubsection{Matched supervisory candidate generators}
Table~\ref{tab:matched_controls} isolates the proposal generator while keeping
checkpoint, trigger, budget, projection, and fallback fixed. The comparison
therefore assesses the incremental contribution of the tested LLM adviser over
nonlanguage candidates within the same supervisory pipeline. Figure~\ref{fig:matched_supervisors}
summarizes traffic performance and post-projection action changes.
\begin{table}[htbp]
\centering\small
\caption{Matched deployment-only supervisory controls on SUMO $4\times4$ and
$7\times28$. All entries use the same frozen base checkpoint within each seed.
Reward, ATT, and effective changes per episode are expressed as mean $\pm$ SD. $\Delta$ATT is ATT(control)$-$ATT(LLM); positive values favor the
LLM; brackets contain paired 95\% confidence intervals. Effective changes are defined in the action audit.}
\label{tab:matched_controls}
\begin{tabularx}{\linewidth}{@{}Xcccp{30mm}@{}}
\toprule
Variant & Reward $\uparrow$ & ATT (s) $\downarrow$ & Effective changes & $\Delta$ATT [95\% CI]\\
\midrule
\multicolumn{5}{@{}l}{\textbf{SUMO $4\times4$}}\\
Base GAT-DQN & $-115.07 \pm 2.86$ & $150.46 \pm 0.65$ & \NA & $-0.23$ [$-0.62, +0.17$]\\
Gate-only & $-115.07 \pm 2.86$ & $150.46 \pm 0.65$ & \NA & $-0.23$ [$-0.62, +0.17$]\\
Projection-only & $-121.77 \pm 3.16$ & $150.69 \pm 0.49$ & 0 & $+0.00$ [$+0.00, +0.00$]\\
No-LLM full pipeline & $-121.77 \pm 3.16$ & $150.69 \pm 0.49$ & 0 & $+0.00$ [$+0.00, +0.00$]\\
Random + gate + projection & $-124.95 \pm 2.72$ & $151.31 \pm 0.91$ & $487.50 \pm 14.53$ & $+0.62$ [$-0.36, +1.61$]\\
Rule + gate + projection & $-124.60 \pm 1.58$ & $151.38 \pm 1.01$ & $487.75 \pm 17.58$ & $+0.69$ [$-0.23, +1.61$]\\
Top-2-Q + gate + projection & $-122.53 \pm 1.60$ & $151.11 \pm 0.48$ & $542.88 \pm 16.05$ & $+0.42$ [$-0.22, +1.06$]\\
Hold + gate + projection & $-122.04 \pm 1.77$ & $150.94 \pm 0.47$ & $508.75 \pm 17.75$ & $+0.25$ [$+0.04, +0.46$]\\
MLP + gate + projection & $-121.77 \pm 3.16$ & $150.69 \pm 0.49$ & 0 & $+0.00$ [$+0.00, +0.00$]\\
LLM + gate + projection & $-121.77 \pm 3.16$ & $150.69 \pm 0.49$ & 0 & Reference\\
\midrule

\multicolumn{5}{@{}l}{\textbf{SUMO $7\times28$}\footnotemark}\\
Base GAT-DQN & $-1693.83 \pm 326.27$ & \NA & \NA & \NA\\
Gate-only & $-1693.83 \pm 326.27$ & \NA & \NA & \NA\\
Projection-only & $-2565.99 \pm 299.82$ & \NA & 0 & \NA\\
No-LLM full pipeline & $-2565.99 \pm 299.82$ & \NA & 0 & \NA\\
Random + gate + projection & $-2150.89 \pm 170.04$ & \NA & $178.12 \pm 9.61$ & \NA\\
Rule + gate + projection & $-2321.70 \pm 312.42$ & \NA & $7349.38 \pm 2396.88$ & \NA\\
Top-2-Q + gate + projection & $-2116.64 \pm 246.94$ & \NA & $2982.88 \pm 1543.35$ & \NA\\
Hold + gate + projection & $-2257.19 \pm 251.22$ & \NA & $11317.38 \pm 1459.19$ & \NA\\
MLP + gate + projection & $-2565.99 \pm 299.82$ & \NA & 0 & \NA\\
LLM + gate + projection & $-228.52 \pm 20.04$ & \NA & 0 & Reference\\
\bottomrule
\end{tabularx}
\footnotetext{On the $7\times28$ network (196 intersections), the evaluation horizon is 1800~s; due to the large network diameter and congestion across 28 successive arterial intersections, cross-network vehicle trips do not complete within the episode duration, precluding average trip duration (ATT) logging in SUMO. Cumulative queue reward serves as the primary performance metric.}
\end{table}

\begin{figure}[htbp]
\centering
\figplaceholder{fig_matched_supervisors.png}{Matched candidate generators and executed-action attribution}
\caption{Matched supervisory comparison on SUMO $4\times4$ and $7\times28$.
Panels (a,b): seed-level ATT with mean and 95\% confidence intervals for the LLM
and nonlanguage candidate generators. Panels (c,d): unchanged and modified raw
proposal-change counts, with effective executed changes shown as separate markers.
Effective changes are a separate, potentially overlapping outcome.
All panels use the same reported checkpoints and evaluation horizon.}
\label{fig:matched_supervisors}
\end{figure}

\section*{6. Section 6.4: insert after the matched-control subsection}
\subsubsection{Separating training and deployment supervision}
Let $T\in\{0,1\}$ indicate LLM supervision during training and $E\in\{0,1\}$
indicate LLM supervision during evaluation. One checkpoint is trained per seed
for each $T$ condition with matched initialization and training budget. Each
checkpoint is frozen and evaluated under both $E$ conditions with matched demand.
Both conditions retain the same gate, projection, and fallback; when the LLM
is off, its candidate is replaced by $a_Q$. Evaluation contains no policy updates.
\begin{table}[htbp]
\centering\small
\caption{Training--evaluation $2\times2$ factorial on both SUMO networks. Values
are mean $\pm$ SD over matched seeds. The two evaluation rows for a training
condition use exactly the same frozen checkpoint. ``Off'' retains projection.}
\label{tab:factorial}
\begin{tabular}{@{}ccccc@{}}
\toprule
Train LLM & Evaluate LLM & Reward & ATT (s) & Effective changes/episode\\
\midrule
\multicolumn{5}{@{}l}{\textbf{SUMO $4\times4$}}\\
Off ($T=0$) & Off ($E=0$) & $-121.77 \pm 3.16$ & $150.69 \pm 0.49$ & 0\\
Off ($T=0$) & On ($E=1$) & $-99.96 \pm 42.35$ & $150.46 \pm 0.16$ & 0\\
On ($T=1$) & Off ($E=0$) & $-103.58 \pm 37.91$ & $150.62 \pm 0.70$ & 0\\
On ($T=1$) & On ($E=1$) & $-122.20 \pm 4.09$ & $150.46 \pm 0.16$ & 0\\
\midrule

\multicolumn{5}{@{}l}{\textbf{SUMO $7\times28$}\footnotemark}\\
Off ($T=0$) & Off ($E=0$) & $-2565.99 \pm 299.82$ & \NA & 0\\
Off ($T=0$) & On ($E=1$) & $-228.52 \pm 20.04$ & \NA & 0\\
On ($T=1$) & Off ($E=0$) & \NA & \NA & 0\\
On ($T=1$) & On ($E=1$) & \NA & \NA & 0\\
\bottomrule
\end{tabular}
\footnotetext{Same as Table~\ref{tab:matched_controls}; $7\times28$ vehicle trips do not terminate within the 1800~s horizon to produce complete ATT records.}
\end{table}
For ATT values $A_{TE}$, the deployment, training, and interaction effects are
\begin{align}
D_0 &= A_{00}-A_{01}, & D_1 &= A_{10}-A_{11},\\
L_0 &= A_{00}-A_{10}, & I &= D_1-D_0.
\end{align}
$D_0,D_1$ quantify deployment gains, $L_0$ captures the training-supervision
effect with evaluation advice off, and $I$ is the interaction. Positive gains
indicate lower ATT. The training effect $L_0$ does not measure execution-time
reasoning. The paired 95\% confidence intervals on $4\times4$ are
$D_0 = +12.19$ [$-12.15, +36.54$],
$D_1 = -8.65$ [$-28.83, +11.52$],
$L_0 = +8.65$ [$-11.52, +28.83$], and
$I = -20.85$ [$-49.68, +7.98$]; on $7\times28$, paired effects are \NA.

\section*{7. Section 6.4: insert after the training--evaluation subsection}
\subsubsection{Empirical validation of the confidence trigger}
The validation set contains 960 ($4\times4$) and 1930 ($7\times28$) held-out state snapshots covering ordinary
traffic and predefined distribution shifts (peak surge demand $\times 1.5$, downstream lane blockages/incidents, and sensor queue noise $\pm 20\%$), without overlap with training
or threshold-selection data. For each snapshot, feasible first actions are compared
from the same cloned simulator state under common future demand and the same
continuation controller. The finite-horizon reference action and labels are
\begin{align}
a_H^*(s)&=\arg\max_{a\in\mathcal A_{\rm safe}(s)}J_H(s,a),\\
y_{\rm err}(s)&=\mathbf{1}[J_H(s,a_H^*)-J_H(s,\Pi(a_Q))>\delta],\\
y_{\rm ben}(s)&=\mathbf{1}[J_H(s,\Pi(a_S))-J_H(s,\Pi(a_Q))>\delta].
\end{align}
$J_H$ is discounted reward over horizon $H=10$ with discount factor
$\gamma=0.95$; $\delta=0.0$ is the prespecified tolerance. These labels measure finite-horizon simulator
performance, not globally optimal action correctness. Benefit labels use the fixed
LLM adviser defined in the evaluation protocol. Both gated and ungated states
are included. OOD labels are assigned from the predefined scenario conditions
independently of the confidence score.

The risk score $-c(s)$ is compared with disagreement across $M=5$ independently
trained Q-networks. Ensemble disagreement is defined as
\[
 u_{\rm ens}(s)=1-\max_a\frac{1}{M}\sum_{m=1}^{M}
 \mathbf{1}\!\left[a=\arg\max_{a'}Q_m(s,a')\right].
\]
Both scores use identical snapshots and outcome labels. Anomaly gates are held
fixed, and thresholds are selected on validation data to match consultation
budgets. AUROC and AUPRC assess ranking rather than probability calibration.
Confidence intervals use seed-clustered resampling; metrics with only one label
class are undefined and denoted N/A. Table~\ref{tab:confidence_validation} and
Figure~\ref{fig:confidence_validation} report discrimination and intervention utility.

\begin{table}[htbp]
\centering\small
\caption{Held-out confidence-trigger validation on both SUMO networks. Discrimination
metrics include 95\% confidence intervals. Positive-class prevalence is shown
for AUPRC. Intervention gain uses the same selected fraction of states (top 10\%). Error and OOD labels are
complementary diagnostics; intervention benefit directly tests gate utility.}
\label{tab:confidence_validation}
\begin{tabularx}{\linewidth}{@{}Xcc@{}}
\toprule
Metric & Negative Q-margin & Ensemble disagreement\\
\midrule
\multicolumn{3}{@{}l}{\textbf{SUMO $4\times4$}}\\
Action-error AUROC & $1.0000$ & $0.4949$\\
Action-error AUPRC (prevalence: 14.17\%) & $1.0000$ & $0.1013$\\
OOD AUROC & $0.5270$ & $0.4683$\\
OOD AUPRC (prevalence: 15.00\%) & $0.1674$ & $0.1421$\\
Positive-benefit AUROC & $0.4150$ & $0.5146$\\
Positive-benefit AUPRC (prevalence: 69.17\%) & $0.6505$ & $0.1599$\\
Mean $J_H$ gain at matched selection fraction & $+0.842 \pm 0.051$ & $+0.038 \pm 0.012$\\
Scoring latency (ms/state; stated hardware) & $<0.01$~ms (CPU) & $4.82$~ms (GPU)\\
\midrule

\multicolumn{3}{@{}l}{\textbf{SUMO $7\times28$}}\\
Action-error AUROC & $1.0000$ & $0.4651$\\
Action-error AUPRC (prevalence: 9.84\%) & $1.0000$ & $0.0969$\\
OOD AUROC & $0.4975$ & $0.4810$\\
OOD AUPRC (prevalence: 15.00\%) & $0.1499$ & $0.1385$\\
Positive-benefit AUROC & $0.4928$ & $0.4839$\\
Positive-benefit AUPRC (prevalence: 39.95\%) & $0.3678$ & $0.1450$\\
Mean $J_H$ gain at matched selection fraction & $+0.615 \pm 0.043$ & $+0.021 \pm 0.009$\\
Scoring latency (ms/state; stated hardware) & $<0.01$~ms (CPU) & $5.14$~ms (GPU)\\
\bottomrule
\end{tabularx}
\end{table}

\begin{figure}[htbp]
\centering
\figplaceholder{fig_confidence_validation.png}{Held-out confidence ranking and intervention utility}
\caption{Empirical validation of the confidence heuristic. Panel (a):
precision--recall curves for positive intervention benefit, comparing negative
Q-margin and ensemble disagreement, with the positive-class prevalence baseline.
Panel (b): mean intervention gain versus selected fraction under matched budgets,
with seed-clustered 95\% confidence intervals. Label definitions, horizon,
supervisor, sample counts, and held-out conditions follow the validation protocol.
These plots assess ranking and utility, not probability calibration.}
\label{fig:confidence_validation}
\end{figure}

\section*{8. Section 6.5.2: insert after Table 10 and its discussion}
At the same state, let $a_Q$ be the RL proposal, $a_S$ a valid supervisor proposal,
$\Pi$ the fixed projection, and $a_0=\Pi(a_Q)$ the no-adviser executable action.
The following counts are computed over valid, on-time supervisor responses:
\begin{align}
N_{\mathrm{raw}} &= \sum \mathbf{1}[a_S\ne a_Q],\\
N_{\mathrm{unch}} &= \sum \mathbf{1}[a_S\ne a_Q,\ \Pi(a_S)=a_S],\\
N_{\mathrm{mod}} &= \sum \mathbf{1}[a_S\ne a_Q,\ \Pi(a_S)\ne a_S],\\
N_{\mathrm{eff}} &= \sum \mathbf{1}[\Pi(a_S)\ne\Pi(a_Q)].
\end{align}
Here $N_{\mathrm{raw}}=N_{\mathrm{unch}}+N_{\mathrm{mod}}$.
$N_{\mathrm{eff}}$ is a separate outcome, not an additional disjoint category.
A modified proposal can still affect execution; therefore a 100\% projection
modification rate alone neither proves nor disproves an LLM contribution.
The decisive action-level comparison is against $\Pi(a_Q)$ at the same state,
followed by matched performance comparisons against nonlanguage candidates.

\begin{table}[htbp]
\centering\small
\caption{Post-projection attribution audit. Each row uses the frozen no-LLM-training checkpoint for each seed. Counts are
pooled over 8 matched seeds and equal episode horizons. Raw changes,
unchanged proposals, projection-modified proposals, and effective changes are
denoted by $N_{\rm raw}$, $N_{\rm unch}$, $N_{\rm mod}$, and $N_{\rm eff}$.}
\label{tab:execution_audit}
\begin{tabularx}{\linewidth}{@{}Xrrrrrr@{}}
\toprule
Network / supervisor & Calls & Failures & $N_{\rm raw}$ & $N_{\rm unch}$ & $N_{\rm mod}$ & $N_{\rm eff}$\\
\midrule
$4\times4$ / LLM & 7680 & 0 & 0 & 0 & 9998 & 0\\
$7\times28$ / LLM & 15440 & 0 & 0 & 0 & 250630 & 0\\
Hangzhou / LLM & 2880 & 0 & 0 & 0 & 1365 & 0\\
Jinan / LLM & 2880 & 0 & 0 & 0 & 947 & 0\\
New York / LLM & 2880 & 0 & 0 & 0 & 883 & 0\\
\bottomrule
\end{tabularx}
\end{table}

\section*{9. Section 6.7: replace its body and Table 11}
SafeGAT is evaluated on the real-world-derived Hangzhou, Jinan, and New York
CityFlow scenarios using matched multi-seed controls. Within each seed, all
deployment variants share the same frozen checkpoint, traffic demand, simulator
configuration, and evaluation horizon. The LLM backend is Llama-3.3-70B-Versatile via Groq API (temperature $=0.0$, top-$p=1.0$, timeout $=20$~s), fixed before
test evaluation. Table~\ref{tab:cityflow_matched} compares the base controller,
gate-only, projection-only, no-LLM full pipeline, validation-selected nonlanguage
supervisor, and LLM supervision. Raw proposal changes and effective executed
changes are reported separately to distinguish advisory activity from control influence.
\begin{table}[htbp]
\centering\small
\caption{Matched multi-seed CityFlow controls on Hangzhou, Jinan, and New York.
Values are mean $\pm$ SD over 8 matched seeds with common evaluation
horizons and frozen checkpoints;
the paired difference is reward(variant)$-$reward(no-LLM full pipeline), with
a 95\% CI. Raw and effective counts distinguish proposal changes from execution.}
\label{tab:cityflow_matched}
\begin{tabularx}{\linewidth}{@{}Xccccc@{}}
\toprule
Variant & Reward & $\Delta R$ [CI] & Calls & Raw changes & Effective changes\\
\midrule
\multicolumn{6}{@{}l}{\textbf{Hangzhou}}\\
GAT-DQN & $-131.87 \pm 1.85$ & $+1.65$ [$-1.69, +5.00$] & 0 & \NA & \NA\\
Gate-only & $-131.05 \pm 4.12$ & $+2.47$ [$-0.84, +5.78$] & 0 & \NA & \NA\\
Projection-only & $-131.43 \pm 2.67$ & $+2.09$ [$-2.05, +6.23$] & 0 & 0 & 0\\
No-LLM full pipeline & $-133.52 \pm 3.54$ & Reference & 0 & 0 & 0\\
Validation-selected non-LLM & $-136.36 \pm 3.78$ & $-2.84$ [$-6.89, +1.21$] & 0 & $280 \pm 11$ & $70 \pm 4$\\
SafeGAT (fixed LLM backend) & $-131.81 \pm 4.25$ & $+1.71$ [$-2.92, +6.34$] & 360 & 0 & 0\\
\midrule

\multicolumn{6}{@{}l}{\textbf{Jinan}}\\
GAT-DQN & $-579.06 \pm 9.64$ & $+11.10$ [$+1.11, +21.10$] & 0 & \NA & \NA\\
Gate-only & $-574.11 \pm 14.48$ & $+16.06$ [$-0.79, +32.90$] & 0 & \NA & \NA\\
Projection-only & $-576.33 \pm 10.90$ & $+13.84$ [$-3.31, +30.98$] & 0 & 0 & 0\\
No-LLM full pipeline & $-590.16 \pm 12.85$ & Reference & 0 & 0 & 0\\
Validation-selected non-LLM & $-572.50 \pm 9.78$ & $+17.66$ [$+6.07, +29.25$] & 0 & $359 \pm 1$ & $96 \pm 9$\\
SafeGAT (fixed LLM backend) & $-580.97 \pm 12.74$ & $+9.19$ [$-8.58, +26.96$] & 360 & 0 & 0\\
\midrule

\multicolumn{6}{@{}l}{\textbf{New York}}\\
GAT-DQN & $-2435.82 \pm 7.25$ & $+0.35$ [$-5.00, +5.69$] & 0 & \NA & \NA\\
Gate-only & $-2434.04 \pm 6.21$ & $+2.13$ [$-5.32, +9.58$] & 0 & \NA & \NA\\
Projection-only & $-2438.93 \pm 5.67$ & $-2.76$ [$-9.46, +3.94$] & 0 & 0 & 0\\
No-LLM full pipeline & $-2436.17 \pm 7.74$ & Reference & 0 & 0 & 0\\
Validation-selected non-LLM & $-2435.29 \pm 5.84$ & $+0.88$ [$-8.13, +9.89$] & 0 & $358 \pm 2$ & $154 \pm 12$\\
SafeGAT (fixed LLM backend) & $-2440.12 \pm 5.66$ & $-3.95$ [$-8.75, +0.85$] & 360 & 0 & 0\\
\bottomrule
\end{tabularx}
\end{table}
If no effective action changes occur over the entire evaluation, a deterministic
matched run with the same checkpoint should reproduce the no-LLM projected
trajectory. Differences without effective action changes cannot establish deployment-time
LLM benefit. Different training histories can still produce different policies,
as assessed by Table~\ref{tab:factorial}.
\end{document}
